% Shared preamble for He (Herb) Huang CV.
% Entry files must set \cvmode to full, short, or onepage before \input{this}.
% Do not compile this file directly.

\RequirePackage[l2tabu,orthodox]{nag}

% Mode flags: extended ⊂ medium ⊂ core
%   full    -> core + medium + extended
%   short   -> core + medium (denser)
%   onepage -> core only (tightest)
\newif\ifcvmedium
\newif\ifcvextended
\ifx\cvmode\undefined
  \errmessage{CV mode not set. Compile cv-hhuang.tex, cv-hhuang.short.tex, or cv-hhuang-onepage.tex}
\fi
\def\cvmodecheckfull{full}
\def\cvmodecheckshort{short}
\def\cvmodecheckonepage{onepage}
\ifx\cvmode\cvmodecheckfull
  \cvmediumtrue
  \cvextendedtrue
\else\ifx\cvmode\cvmodecheckshort
  \cvmediumtrue
  \cvextendedfalse
\else\ifx\cvmode\cvmodecheckonepage
  \cvmediumfalse
  \cvextendedfalse
\else
  \errmessage{Unknown \string\cvmode; use full, short, or onepage}
\fi\fi\fi

% Encoding and production typography (pdflatex).
%
% EB Garamond carries the body and the name (tabular lining figures: years
% and dates sit on the same baseline as the rest of the line). Alegreya Sans
% carries the three heading levels, using the family's native Bold and
% Medium Italic faces. Keep these defaults explicit: both packages set NFSS
% family defaults while loading, and package order should not silently
% change the finished CV.
\usepackage[T1]{fontenc}
\usepackage[utf8]{inputenc}
\usepackage[lining,tabular]{ebgaramond}
\usepackage{AlegreyaSans}
\renewcommand*{\rmdefault}{EBGaramond-TLF}
\renewcommand*{\sfdefault}{AlegreyaSans-LF}
\newcommand{\cvheadfamily}{\sffamily}
\newcommand{\cvheaddefault}{\sfdefault}

\usepackage[USenglish]{babel}
\usepackage[babel=true]{microtype}

\usepackage{datetime}
\usepackage{geometry}
\usepackage{setspace}
\usepackage{titlesec}
\usepackage[super]{nth}
\usepackage{hyperref}

\newcommand{\myname}{He (Herb) Huang}

% Density by mode (margins only — avoid body height vs top/bottom conflicts).
% Leading is 1.0 inside a paragraph and \listitemspace adds the extra 0.20 of
% a line between entries, so entry-to-entry spacing reads as 1.20. The short
% and one-page modes trade that away -- they carry the same sections as the
% full CV and only fit their page budget compressed.

% Left label column + hanging indent for wrapped entry lines. The same in every
% mode: it aligns columns down the page, so it is not a density knob.
\newcommand{\listtabwidth}{0.5in}

\ifcvextended
  % full
  \newcommand{\listitemspace}{0.2\baselineskip}
  \geometry{left=1.0in, right=1.0in, top=1.0in, bottom=1.0in, footskip=24pt}
  \setstretch{1.0}
  \titlespacing{\section}{0pt}{24pt plus 4pt minus 4pt}{9pt plus 2pt minus 2pt}
  \titlespacing{\subsection}{0pt}{8pt plus 4pt minus 4pt}{4pt plus 2pt minus 2pt}
  \titlespacing{\subsubsection}{0pt}{5pt plus 2pt minus 2pt}{2pt plus 1pt minus 1pt}
\else\ifcvmedium
  % short (dense multi-section)
  \newcommand{\listitemspace}{0.04\baselineskip}
  \geometry{left=0.75in, right=0.75in, top=0.6in, bottom=0.6in, footskip=18pt}
  \setstretch{0.90}
  \titlespacing{\section}{0pt}{12pt plus 2pt minus 2pt}{5.5pt plus 1pt minus 1pt}
  \titlespacing{\subsection}{0pt}{5pt plus 2pt minus 2pt}{2pt plus 1pt minus 1pt}
  \titlespacing{\subsubsection}{0pt}{3pt plus 1pt minus 1pt}{1pt plus 0.5pt minus 0.5pt}
\else
  % onepage
  \newcommand{\listitemspace}{0.06\baselineskip}
  \geometry{left=0.5in, right=0.5in, top=0.4in, bottom=0.4in, footskip=14pt}
  \setstretch{0.88}
  \titlespacing{\section}{0pt}{5pt plus 1pt minus 1pt}{2.5pt plus 1pt minus 1pt}
  \titlespacing{\subsection}{0pt}{2.5pt plus 1pt minus 1pt}{1pt plus 0.5pt minus 0.5pt}
  \titlespacing{\subsubsection}{0pt}{2pt plus 1pt minus 1pt}{1pt plus 0.5pt minus 0.5pt}
\fi\fi

\makeatletter
\newcommand{\namefont}[1]{{\normalfont\rmfamily\bfseries\Large
  \fontsize{\dimexpr\f@size pt+1pt\relax}{\dimexpr\f@size pt+4pt\relax}%
  \selectfont #1}}
\makeatother

% Own-name emphasis in reference lists.
\newcommand{\me}[1]{\textbf{#1}}

\SetTracking{encoding=*, family=\cvheaddefault}{30}
% Rule under level-1 headings, overhanging the measure by \cvruleover at each
% end. It goes through titlesec's \titleline so it stays a bare box on the
% vertical list -- a \makebox here would open a paragraph and pick up a full
% \baselineskip of interline glue above and below the rule. The zero-width
% \makebox lets the rule overhang without reporting an overfull box.
\newlength{\cvruleover} \setlength{\cvruleover}{2pt}
\newcommand{\cvsectionrule}{%
  \vspace{-2pt}%
  \titleline{%
    \makebox[0pt][l]{\hspace*{-\cvruleover}%
      \rule{\dimexpr\linewidth+2\cvruleover\relax}{0.4pt}}\hfil}}
% Level 1: Alegreya Sans Bold, uppercase and tracked.
\titleformat{\section}{\lsstyle\cvheadfamily\normalsize\bfseries\uppercase}{}{}{}%
  [{\cvsectionrule}]
% Level 2: Alegreya Sans Bold, title case and lightly tracked.
\titleformat{\subsection}{\lsstyle\cvheadfamily\normalsize\bfseries}{}{}{}{}
% Level 3: Alegreya Sans Medium Italic -- distinct without competing with L2.
\titleformat{\subsubsection}%
  {\cvheadfamily\normalsize\fontseries{sb}\itshape\selectfont}{}{0pt}{}

\setlength\parindent{0em}

% Uniform word spacing. TeX otherwise pads the space after `.', `?' and `:'
% to ~1.5x, which in a reference list lands unevenly -- an entry whose title
% ends in a sentence gets loose lines while an all-initials author string
% stays tight. APA reference lists are single-spaced after periods anyway.
%\frenchspacing

% ---------------------------------------------------------------------------
% Entry layout
%
%   \entl{PhD}{2021--present}  ->  PhD   body text ...............  2021--present
%                                        wrapped lines hang here
%   \ent{2025}                 ->  body text .............................  2025
%                                      wrapped lines hang here
%   \entn                      ->  body text, full width, no date
%                                      wrapped lines hang here
%
% All three hang wrapped lines at \cvtabw so every column lines up down the
% page. Use them inside `entries`, which neutralises the \par that \raggedright
% smuggles into \\ (a \par would reset \parshape and drop the hang).
% ---------------------------------------------------------------------------
% The date column is measured at \begin{document} from the widest notation the
% document actually prints, so swapping the body font resizes it automatically.
% \cvdatewwide is the variant for entries that carry two terms at once; pass it
% as \begin{entries}[\cvdatewwide].
\newcommand{\cvdatesample}{2021--present}
\newcommand{\cvdatesamplewide}{Fall 2022; Spring 2025}
\newlength{\cvtabw}   \setlength{\cvtabw}{\listtabwidth}
\newlength{\cvgap}    \setlength{\cvgap}{0.6em}
\newlength{\cvminlast}\setlength{\cvminlast}{6em}
\newlength{\cvragged} \setlength{\cvragged}{3em}
\newlength{\cvdatew}
\newlength{\cvdatewwide}
\newlength{\cvmeas}
\AtBeginDocument{%
  \settowidth{\cvdatew}{\cvdatesample}\addtolength{\cvdatew}{0.35em}%
  \settowidth{\cvdatewwide}{\cvdatesamplewide}\addtolength{\cvdatewwide}{0.35em}}

% Zero-width overlay: puts #1 flush against the right margin without consuming
% horizontal space, so the body's own line breaking is untouched.
\newcommand{\cvdate}[1]{%
  \settowidth{\cvmeas}{#1}%
  \ifdim\cvmeas>\cvdatew
    \PackageWarning{cv}{Date '#1' overruns the column
      (\the\cvmeas\space vs \the\cvdatew); widen \string\cvdatesample}%
  \fi
  \makebox[0pt][l]{\hspace*{\dimexpr\linewidth-\cvdatew\relax}%
    \makebox[\cvdatew][r]{#1}}}

% Every line stops short of the date column, so the body forms a rectangular
% block that never runs underneath the date. Lines 2+ hang at \cvtabw.
% \parfillskip caps how far the final line may fall short of the measure,
% which is what forbids a one-word stub at the end of an entry.
\newcommand{\cvshape}{%
  \par\addvspace{\listitemspace}%
  \parshape 2 0pt \dimexpr\linewidth-\cvdatew-\cvgap\relax
              \cvtabw \dimexpr\linewidth-\cvtabw-\cvdatew-\cvgap\relax
  \parfillskip=0pt plus
    \dimexpr\linewidth-\cvtabw-\cvdatew-\cvgap-\cvminlast-\cvragged\relax
  \noindent}

\newcommand{\ent}[1]{\cvshape\cvdate{#1}\ignorespaces}


\newcommand{\entl}[2]{%
  \settowidth{\cvmeas}{#1}%
  \ifdim\cvmeas>\dimexpr\cvtabw-0.5em\relax
    \PackageWarning{cv}{Label '#1' crowds the \string\listtabwidth\space column}%
  \fi
  \cvshape\cvdate{#2}\makebox[\cvtabw][l]{#1}\ignorespaces}

\newcommand{\entn}{%
  \par\addvspace{\listitemspace}%
  \hangindent=\cvtabw \hangafter=1
  \parfillskip=0pt plus \dimexpr\linewidth-\cvtabw-\cvminlast-\cvragged\relax
  \noindent\ignorespaces}

\makeatletter
\newenvironment{entries}[1][\cvdatew]
  {\par\begingroup
   \setlength{\cvdatew}{#1}%
   \setlength{\parskip}{0pt}%
   \let\\\@normalcr
   % \raggedright's \rightskip is 0pt plus 1fil: infinite stretch means every
   % line scores badness 0, so \parfillskip cannot forbid anything. Swapping in
   % a finite \rightskip restores real badness; \cvragged is how far short of
   % the measure an ordinary line may fall, so it doubles as the cap on how much
   % the right edge may jitter -- but it must stay wide enough that the breaker
   % can still afford to push a word down onto the final line. \hyphenpenalty
   % keeps that tightness from being paid for in broken words, and
   % \exhyphenpenalty does the same for the hyphens words already carry
   % (Person-Supervisor), which \hyphenpenalty does not reach.
   \rightskip=0pt plus \cvragged
   \hyphenpenalty=2000 \exhyphenpenalty=2000
   \tolerance=9999 \emergencystretch=2em \hbadness=4000}
  {\par\endgroup}
\makeatother

% Two-column block (address | contact, or reference | contact). The right
% column is set at its natural width so it can never wrap, and the left column
% takes what is left; unlike fixed \textwidth fractions this survives a change
% of body font.
\newsavebox{\cvhdrbox}
\newcommand{\cvheader}[2]{%
  \sbox{\cvhdrbox}{\begin{tabular}[t]{@{}r@{}}#2\end{tabular}}%
  \par\noindent
  \hbox to\textwidth{%
    \begin{minipage}[t]{\dimexpr\textwidth-\wd\cvhdrbox-1em\relax}#1\end{minipage}%
    \hfill\usebox{\cvhdrbox}}%
  \par}

% ---------------------------------------------------------------------------
% Author marks and the trailing note on a publication entry.
%   \eqc / \advisee   superscript marks used inside author lists; \cvmarks is
%                     the key, and lists only the marks in use
%   \jabdc ... \jutd  one per ranking list, each linked to the list's own site
%   \cvnote{...}      the single bracketed group that closes an entry
%
% Everything that trails an entry -- review status and ranking lists alike --
% goes in one \cvnote, semicolon separated, status first:
%
%   \cvnote{2nd~R\&R; \jabs; \jft; \jtamuga}
%     -> [2nd R&R; ABS 4*; FT50; TAMUGA]
%
% The group may break after a `;' -- that reads fine -- so tie any multi-word
% status (2nd~R\&R, Fully~accepted) to keep it off a line of its own.
%
% Memberships verified 2026-08-22: TAMUGA indexes ten management journals
% (JAP yes, JIBS no); JIBS is one of the UTD24.
% ---------------------------------------------------------------------------
\newcommand{\eqc}{\textsuperscript{*}}
\newcommand{\advisee}{\textsuperscript{\dag}}
% The printed key lists only the marks the document actually uses -- a key that
% explains a symbol the reader never meets is noise. Add the advisee clause
% back the first time an entry carries \advisee.
\newcommand{\cvmarks}{\eqc Equal contribution.}

% \jabdc is defined but deliberately not printed: ABDC A* adds nothing next to
% ABS 4* and FT50. Drop it into a \cvnote if you ever want it back.
\newcommand{\jabdc}{\href{https://abdc.edu.au/abdc-journal-quality-list/}%
  {ABDC\,A\textsuperscript{*}}}
\newcommand{\jabs}{\href{https://charteredabs.org/academic-journal-guide}%
  {ABS\,4\textsuperscript{*}}}
\newcommand{\jft}{\href{https://www.ft.com/content/3405a512-5cbb-11e1-8f1f-00144feabdc0}%
  {FT50}}
\newcommand{\jtamuga}{\href{https://www.tamugarankings.com/}{TAMUGA}}
\newcommand{\jutd}{\href{https://jsom.utdallas.edu/the-utd-top-100-business-school-research-rankings/list-of-journals}%
  {UTD24}}
\newcommand{\cvnote}[1]{[#1]}

% The Publications heading carries the mark key: a block underneath it in the
% full CV, inline on the heading itself in the compressed modes where a whole
% line is not worth the space. The inline note is boxed first so the heading's
% \uppercase cannot reach its letters.
\newsavebox{\cvnotebox}
\newcommand{\cvpubheading}{%
  \ifcvextended
    \section*{Publications}%
    \cvlegend{\cvmarks}%
  \else
    \savebox{\cvnotebox}{\normalfont\mdseries\footnotesize\cvmarks}%
    \section*{Publications\hfill\usebox{\cvnotebox}}%
  \fi}

% Section-local mark legend, e.g. the journal-ranking and equal-contribution
% marks used in Publications. Sits directly under the section heading so the
% key is where the reader first meets the marks, rather than at the page foot
% -- which in the full CV is a page away from half the entries that use them.
\newcommand{\cvlegend}[1]{%
  {\par\footnotesize
   % No \addvspace above: the section's own after-spacing already separates the
   % key from the rule, and a second gap detaches it from the heading it keys.
   \rightskip=0pt plus 3em\relax #1\par
   \addvspace{0.35\baselineskip}}}

\newdateformat{monthyeardate}{\monthname[\THEMONTH] \THEYEAR}

\hypersetup{
    colorlinks  = true,
    urlcolor    = black,
    citecolor   = black,
    linkcolor   = black,
    pdfauthor   = {\myname},
    pdfkeywords = {organizational behavior, human resources, management, employee mobility, cognition, cv, academic},
    pdftitle    = {\myname: Curriculum Vitae},
    pdfsubject  = {Curriculum Vitae},
    pdfpagemode = UseNone
}
